% Generated draft: requires mathematical and source review.
\documentclass[11pt]{article}
\usepackage[margin=25mm]{geometry}
\usepackage{fontspec}
\setmainfont{Latin Modern Roman}
\usepackage{amsmath,amssymb,mathtools,unicode-math}
\setmathfont{Latin Modern Math}
\usepackage{enumitem,needspace,xurl,hyperref}
\setlength{\parindent}{0pt}
\setlength{\parskip}{5pt}
\setlength{\emergencystretch}{3em}
\newcommand{\N}{\mathbb{N}}
\newcommand{\Z}{\mathbb{Z}}
\newcommand{\Q}{\mathbb{Q}}
\newcommand{\R}{\mathbb{R}}
\newcommand{\C}{\mathbb{C}}
\newcommand{\F}{\mathbb{F}}
\newcommand{\A}{\mathbb{A}}
\newcommand{\PP}{\mathbb{P}}
\newcommand{\E}{\mathbb{E}}
\newcommand{\eps}{\varepsilon}
\newcommand{\dd}{\,\mathrm{d}}
\begin{document}
\section*{UnsolvedMath ID: 30006513}
\textbf{Source problem code:} catalogue entry
\subsection{Definitions and mathematical statement}
A first-order structure \(M\) is \(\omega\)-categorical if its complete theory has, up to isomorphism, a unique countable model. Put \(G=\operatorname{Aut}(M)\). A \(G\)-set \(W\) is oligomorphic if, for every integer \(n\geq 1\), the diagonal action of \(G\) on \(W^n\) has finitely many orbits.

For a field \(F\), let \(FG\) be the group algebra of the abstract group \(G\), and define the permutation module
\[
F[W]=\bigoplus_{w\in W}F e_w,\qquad g\cdot e_w=e_{g\cdot w}.
\]
Thus its elements are finite linear combinations of basis vectors indexed by \(W\). An \(FG\)-module \(V\) has finite composition length if there are an integer \(\ell\geq 0\) and a chain of \(FG\)-submodules
\[
0=V_0\subsetneq V_1\subsetneq\cdots\subsetneq V_\ell=V
\]
such that every quotient \(V_i/V_{i-1}\), for \(1\leq i\leq\ell\), is nonzero and has no proper nonzero \(FG\)-submodule. The case \(\ell=0\) means \(V=0\).

\textbf{Question as supplied.} For every \(\omega\)-categorical structure \(M\), every oligomorphic \(\operatorname{Aut}(M)\)-set \(W\), and every field \(F\), does the module \(F[W]\) have finite composition length over \(F\operatorname{Aut}(M)\)? There is no restriction on the characteristic of \(F\) in the supplied statement.

\textbf{Unresolved scope qualifications.} The supplied statement does not specify the cardinality of \(M\), the size of its language, or whether the action on \(W\) must be continuous when \(\operatorname{Aut}(M)\) carries its usual permutation topology and \(W\) is discrete. The question above retains the supplied quantifiers; it does not silently add these restrictions. Oligomorphicity alone is not being used as a substitute for an explicit continuity assumption.

The supplied literature assessment discusses the narrower class of homogeneous structures in finite relational languages. That restriction is absent from the supplied mathematical statement, and the incomplete original excerpt does not establish whether it was an intended hypothesis. The excerpt subsequently discusses ordered Ramsey structures, extreme amenability, and a dual module of definable functions; these are presented as surrounding exposition, not as additional hypotheses of the question above. The notation \(F[W]\) here follows the supplied cleaned statement and denotes the finite-support permutation module, not that separate function module.
\subsection{Short English statement}
The recorded question asks whether every oligomorphic action of the automorphism group of an \(\omega\)-categorical structure yields a permutation module with a finite composition series over every field. The intended restrictions on the structure and action require review. The supplied status is ``partially solved,'' based on dated curation rather than independent verification here.
\subsection{Sources}
\begin{itemize}
\item Local UnsolvedMath registry, permanent ID \(30006513\), source problem code \texttt{catalogue entry}, title \emph{Finite Length of Permutation Modules over Arbitrary Fields}; supplied complete source record from the Oberwolfach Reports Open Problems collection. No proposer is identified in the record.
\item Supplied primary provenance: \emph{Homogeneous Structures: Model Theory meets Universal Algebra} (2026), Oberwolfach Reports, DOI \texttt{10.4171/owr/2025/58}, pp.\ 3063--3124. \url{https://doi.org/10.4171/owr/2025/58}. The record gives proposed year 2025; the relationship between this date, the DOI year, and the displayed citation year 2026 is unresolved.
\item Supplied literature source: Evans, \emph{Permutation modules for Ramsey structures}. \url{https://arxiv.org/abs/2603.29606}.
\item Supplied literature source: Yang--Boja\'nczyk--Klin, \emph{The Finite Length Property of the Rado Graph and Friends}. \url{https://www.cs.ox.ac.uk/people/bartek.klin/papers/lics26.pdf}.
\item Literature assessment dated exactly \texttt{2026-08-22}: the supplied curation labels the problem ``partially solved'' and reports finite composition length for the Rado graph and several broad classes, while reporting that the question for every finite-relational homogeneous structure remains open. This is supplied dated curation derived from the OWR record, not an independently verified literature assessment; the cited claims and time-sensitive status require independent checking.
\item The supplied registry labels the statement \texttt{corrected\_verified} and reports that the official Oberwolfach report was checked, the question isolated, and its notation repaired. This is a supplied verification claim, not a claim of source inspection in this draft.
\item The original excerpt mentions a question of Macpherson, Question 2.2.7(4) in its reference [13], a conjecture of Simon Thomas on reducts of finitely homogeneous structures in its reference [14], and a theorem of Kechris, Pestov and Todor\v{c}evi\'c in its reference [12]. The associated bibliography and the antecedents of the excerpt's clauses (i) and (ii) are not supplied, so these references cannot be identified more precisely from this record.
\end{itemize}
\end{document}
